\uuid{z6r2}
\exo7id{6307}
\titre{exo7 6307}
\auteur{mayer}
\organisation{exo7}
\datecreate{2011-10-16}
\isIndication{false}
\isCorrection{false}
\chapitre{Extremum, extremum lié}
\sousChapitre{Extremum, extremum lié}
\cpge{mpsi-mp2i.fonctions-deux-variables.b, pcsi.fonctions-deux-variables.b, ptsi.fonctions-deux-variables.b, mpsi-mp2i.fonctions-deux-variables.d, pcsi.fonctions-deux-variables.d, ptsi.fonctions-deux-variables.d, bcpst1.derivees-primitives-integrales.b, bcpst1.fonctions-deux-variables.c, mp-mpi.calcul-differentiel-optimisation.f}
\module{Calcul différentiel}
\niveau{L3}
\difficulte{}

\contenu{
\question{
Discuter, suivant les valeurs du paramètre réel $\lambda $, la
nature des extremums de la fonction $f(x,y) = y(x^2+y^2 -2\lambda y)$.
}
}
